% =====================================================================
% Preambuła książki naukowej wiernie wzorowana na „Mathematics for Machine
% Learning” (Deisenroth, Faisal, Ong):
%  - czcionka Times, czarno-biały skład, bez kolorowych ramek,
%  - definicje i twierdzenia: pogrubiony nagłówek w tekście, bez ramki,
%  - szare tło TYLKO dla przykładów,
%  - szeroki margines: podpisy rysunków, hasła (\haslo) i uwagi na marginesie,
%  - uwagi i intuicje jako kursywne „Uwaga.” zakończone znakiem ◊,
%  - nagłówek strony: numer + kursywny tytuł rozdziału / sekcji, bez linii.
% Wykresy powstają osobno w Pythonie (figures/) i wchodzą jako obrazy.
% =====================================================================
\usepackage[T1]{fontenc}
\usepackage[utf8]{inputenc}
\usepackage[polish]{babel}
\usepackage{mathptmx}              % Times + dopasowana matematyka
\usepackage{microtype}
\usepackage{amsmath,amssymb,amsthm,mathtools,bm}
\usepackage{booktabs,array}
\usepackage{graphicx}
\usepackage{xcolor}
\usepackage{enumitem}
\setlist[enumerate]{leftmargin=*,itemsep=2pt,topsep=3pt}
\setlist[itemize]{leftmargin=*,itemsep=2pt,topsep=3pt}

% ---------- układ strony: szeroki zewnętrzny margines ----------
\usepackage[a4paper,twoside,inner=2.4cm,textwidth=11.6cm,marginparsep=0.55cm,
  marginparwidth=4.9cm,top=2.7cm,bottom=2.9cm,headheight=14pt,headsep=0.9cm]{geometry}
\usepackage{sidenotes}             % \sidenote, \sidecaption, marginfigure, figure*
\definecolor{mmlgray}{gray}{0.94}  % jedyny „kolor” książki: szare tło przykładów

\usepackage{hyperref}
\hypersetup{colorlinks=true,linkcolor=black,citecolor=black,urlcolor=black,
  pdfauthor={},pdftitle={}}
\usepackage{cleveref}
\crefname{figure}{rys.}{rys.}\Crefname{figure}{Rysunek}{Rysunki}
\crefname{table}{tab.}{tab.}\Crefname{table}{Tabela}{Tabele}
\crefname{equation}{wzór}{wzory}\Crefname{equation}{Wzór}{Wzory}
\crefname{chapter}{rozdział}{rozdziały}\Crefname{chapter}{Rozdział}{Rozdziały}
\crefname{section}{sekcja}{sekcje}\Crefname{section}{Sekcja}{Sekcje}
\crefname{definicja}{def.}{def.}\Crefname{definicja}{Definicja}{Definicje}
\crefname{twierdzenie}{tw.}{tw.}\Crefname{twierdzenie}{Twierdzenie}{Twierdzenia}
\crefname{lemat}{lemat}{lematy}\Crefname{lemat}{Lemat}{Lematy}
\crefname{wniosek}{wniosek}{wnioski}\Crefname{wniosek}{Wniosek}{Wnioski}
\crefname{przyklad}{przykład}{przykłady}\Crefname{przyklad}{Przykład}{Przykłady}

\numberwithin{equation}{chapter}   % (2.94): numer rozdziału i wzoru

% ---------- nagłówki rozdziałów i sekcji (jak w MML: czarne, pogrubione) ----------
\usepackage[explicit]{titlesec}
\titleformat{\chapter}[hang]{\bfseries\fontsize{26}{30}\selectfont}{\thechapter}{0.9em}{#1}
\titleformat{name=\chapter,numberless}[hang]{\bfseries\fontsize{26}{30}\selectfont}{}{0pt}{#1}
\titlespacing*{\chapter}{0pt}{16pt}{34pt}
\titleformat{\section}{\large\bfseries}{\thesection}{0.8em}{#1}
\titleformat{\subsection}{\bfseries\itshape}{\thesubsection}{0.8em}{#1}
\titleformat{\paragraph}[runin]{\bfseries\itshape}{}{0pt}{#1.}
\titlespacing*{\section}{0pt}{2.8ex plus 1ex minus .2ex}{1.2ex plus .2ex}
\titlespacing*{\subsection}{0pt}{2.4ex plus 1ex minus .2ex}{1ex plus .2ex}

% ---------- nagłówki i stopki stron (numer + kursywny tytuł, bez linii) ----------
\usepackage{fancyhdr}
\pagestyle{fancy}
\fancyhf{}
\renewcommand{\headrulewidth}{0pt}
\renewcommand{\chaptermark}[1]{\markboth{#1}{}}
\renewcommand{\sectionmark}[1]{\markright{\thesection\ #1}}
\fancyhead[LE]{\thepage}
\fancyhead[RE]{\textit{\leftmark}}
\fancyhead[LO]{\textit{\rightmark}}
\fancyhead[RO]{\thepage}
\fancypagestyle{plain}{\fancyhf{}\fancyfoot[C]{\thepage}\renewcommand{\headrulewidth}{0pt}}

% ---------- spis treści: pogrubione rozdziały, bez kropek ----------
\usepackage{tocloft}
\renewcommand{\cfttoctitlefont}{\bfseries\fontsize{26}{30}\selectfont}
\renewcommand{\cftaftertoctitleskip}{24pt}
\renewcommand{\cftchapfont}{\bfseries}
\renewcommand{\cftchappagefont}{\bfseries}
\renewcommand{\cftdotsep}{\cftnodots}

% ---------- podpisy rysunków: mała czcionka, pogrubiona etykieta, w marginesie ----------
\captionsetup{font=footnotesize,labelfont=bf,justification=raggedright,
  singlelinecheck=off,skip=4pt}

% ---------- twierdzenia, definicje, przykłady ----------
\usepackage[most]{tcolorbox}
% nagłówek pogrubiony w tekście, bez ramki: „Definicja 2.10 (Nazwa).”
\newtheoremstyle{mmldef}{9pt}{9pt}{\normalfont}{0pt}{\bfseries}{.}{0.5em}
  {\thmname{#1}\ \thmnumber{#2}\thmnote{\ (#3)}}
\newtheoremstyle{mmlthm}{9pt}{9pt}{\itshape}{0pt}{\bfseries}{.}{0.5em}
  {\thmname{#1}\ \thmnumber{#2}\thmnote{\ (#3)}}
\theoremstyle{mmldef}
\newtheorem{definicja}{Definicja}[chapter]
\newtheorem{przyklad}{Przykład}[chapter]       % własna numeracja, jak w MML
\theoremstyle{mmlthm}
\newtheorem{twierdzenie}[definicja]{Twierdzenie}
\newtheorem{lemat}[definicja]{Lemat}
\newtheorem{wniosek}[definicja]{Wniosek}
\renewcommand{\qedsymbol}{$\blacksquare$}

% jedyna ramka: szare tło przykładów
\tcbset{mmlbox/.style={enhanced,breakable,sharp corners,frame hidden,boxrule=0pt,
  colback=mmlgray,left=8pt,right=8pt,top=5pt,bottom=5pt,
  before skip=11pt,after skip=11pt,before upper={\parindent0pt}}}
\tcolorboxenvironment{przyklad}{mmlbox}
\newtcolorbox{naliczbach}{mmlbox,
  before upper={\parindent0pt\textbf{Na konkretnych liczbach.}\ }}

% uwagi jak w MML: kursywna etykieta w tekście i znak ◊ na końcu
\newenvironment{mmluwaga}[1]{\par\smallskip\noindent\textit{#1.}\ }{\hfill$\diamondsuit$\par\smallskip}
\newenvironment{uwaga}{\begin{mmluwaga}{Uwaga}}{\end{mmluwaga}}
\newenvironment{intuicja}{\begin{mmluwaga}{Intuicja}}{\end{mmluwaga}}
\newenvironment{pulapka}{\begin{mmluwaga}{Pułapka}}{\end{mmluwaga}}
\newenvironment{pomyslDowodu}{\par\smallskip\noindent\textit{Pomysł dowodu.}\ }{\par\smallskip}

% ćwiczenia na końcu rozdziału: zwykła sekcja z listą
\newenvironment{sprawdz}{\subsection*{Ćwiczenia}}{}

% hasło w marginesie (jak terminy w MML): \haslo{wyznacznik}
\newcommand{\haslo}[1]{\marginnote{\footnotesize\raggedright #1}}

% Odesłanie do dowodu w dodatku (dla dowodów zbyt długich na rozdział).
\newcommand{\dowodw}[1]{\par\noindent\textit{Dowód: patrz \cref{#1}.}\par}

% ---------- skróty matematyczne (wektory pogrubione kursywą jak w MML) ----------
\let\Im\relax % w LaTeX-u \Im to część urojona; tu znaczy obraz przekształcenia
\DeclareMathOperator{\Im}{Im}
\DeclareMathOperator{\rk}{rk}
\DeclareMathOperator{\Ker}{Ker}
\DeclareMathOperator{\tr}{tr}
\DeclareMathOperator{\spn}{span}
\DeclareMathOperator*{\argmin}{arg\,min}
\DeclareMathOperator*{\argmax}{arg\,max}
\newcommand{\R}{\mathbb{R}}
\newcommand{\C}{\mathbb{C}}
\newcommand{\N}{\mathbb{N}}
\newcommand{\vect}[1]{\bm{#1}}
\newcommand{\norm}[1]{\left\lVert #1\right\rVert}
\newcommand{\ip}[2]{\left\langle #1,#2\right\rangle}
\newcommand{\abs}[1]{\left\lvert #1\right\rvert}

% ---------- schematy w TikZ (tylko proste diagramy; wykresy robimy w Pythonie) ----------
\usepackage{tikz}
\usetikzlibrary{arrows.meta,calc,positioning,shapes.geometric,fit,backgrounds}
\tikzset{
  wektor/.style={-{Stealth[length=2.2mm]},thick},
  pomocnicza/.style={dashed,gray},
  wezel/.style={circle,draw,thick,minimum size=7mm,inner sep=1pt},
  blok/.style={rectangle,draw,thick,minimum height=8mm,minimum width=18mm,align=center,fill=white}
}

% ---------- strona tytułowa (czarno-biała) ----------
% \tytulksiazki{Tytuł}{Podtytuł}{Autor}
\newcommand{\tytulksiazki}[3]{%
  \begin{titlepage}
  \thispagestyle{empty}
  \vspace*{4.5cm}
  {\raggedright\bfseries\fontsize{40}{46}\selectfont #1\par}
  \vspace{1.0cm}
  {\raggedright\Large #2\par}
  \vfill
  {\raggedright\Large #3\par}
  \vspace{1.5cm}
  \end{titlepage}}

% ---------- makra tej książki (model Gen 6) ----------
\DeclareMathOperator{\softmax}{softmax}
\DeclareMathOperator{\GELU}{GELU}
\DeclareMathOperator{\LN}{LN}
\DeclareMathOperator{\Drop}{Drop}
\DeclareMathOperator{\MLP}{MLP}
\DeclareMathOperator{\Conv}{Conv}
\DeclareMathOperator{\DW}{DW}
\DeclareMathOperator{\shift}{shift}
\DeclareMathOperator{\clip}{clip}
\DeclareMathOperator{\round}{round}
\DeclareMathOperator{\ReLU}{ReLU}
\DeclareMathOperator{\Var}{Var}
\DeclareMathOperator{\Exp}{\mathbb{E}}
\DeclareMathOperator{\sg}{sg}
\newcommand{\indic}[1]{\mathbb{1}\!\left[#1\right]}
\newcommand{\ie}{tzn.\ }
\newcommand{\kodlink}[1]{\texttt{#1}}
\newcommand{\Hz}{\,\mathrm{Hz}}
\newcommand{\dB}{\,\mathrm{dB}}
\newcommand{\ms}{\,\mathrm{ms}}
\newcommand{\pow}[1]{\text{\footnotesize #1}}   % krótka adnotacja przy kroku dowodu
